\section{Current primary inputs on CM Hodge, abelian Tate, and the
standard
conjectures}\label{current-primary-inputs-on-cm-hodge-abelian-tate-and-the-standard-conjectures}

Research module \texttt{consequences\_criteria}, checked 20--21 November
2026 UTC (literature cutoff: 20 November), GPT-6 Astra Pro; effort not
exposed. This module updates the archive analysis using identified 2026
primary manuscripts. Their major theorems are \textbf{cited literature
inputs}, not theorems independently reproved in this run. The adopted
Alpöge RH and Dirichlet GRH hypotheses play no role in the deductions
below.

\subsection{CC-TATE01. The current all-CM input and its exact
provenance}\label{cc-tate01.-the-current-all-cm-input-and-its-exact-provenance}

\textbf{Cited 2026 preprint theorem.} For every complex abelian variety
of CM type, every rational Hodge class in every codimension is
algebraic. The precise locator is van Hoften and Sempliner, \emph{The
Hodge Conjecture for CM abelian varieties}, \textbf{Theorem I, p.~1},
Hexagon \textbf{2610.00259v2}. This includes powers and products of CM
abelian varieties, because those varieties are themselves of CM type.

The version history dates v1 to 15 October and v2 to \textbf{18 October
2026}. The current record's header displays the original 15 October
date; the downloaded v2 PDF creation and modification timestamps both
read 18 October 2026, 10:11:24 UTC. Thus the edition used here is
specified by the immutable v2 identifier, not merely by the header date.
The named authors are Pol van Hoften and Jack Sempliner. Their §1.3,
p.~4, explains that they give a variant proof inspired by the OpenAI
manuscript and do not certify every detail of that earlier manuscript.
Our cited input is their own Theorem I. Neither a journal publication
nor a complete independent verification of their proof is claimed here.

\textbf{Dependency inspection.} The reduction is Theorem 2.2.1, p.~6;
the algebraic surface criterion is Theorem 2.3.4, pp.~7--8; the required
nonzero cup product is Theorem 4.0.2, p.~12. We read the
surface-criterion proof completely, the cited Albanese statement Theorem
3.1.11 and Corollary 3.1.13, pp.~11--12, and the final argument §4.2.8,
pp.~19--22, including the diagonal-period Lemma 4.2.7, pp.~17--18, and
the archimedean nonvanishing Lemma 4.2.9, pp.~20--21. Those analytic
lemmas retain their cited theta-correspondence inputs; these inputs were
not independently reconstructed. The reduction's reference to Hazama is
recorded as a dependency, including the authors' p.~6 footnote about the
formulation of Hazama's Lemmas 7.11--7.12. This is a dependency read,
not a claim to have checked every underlying paper.

Milne's \textbf{9 November 2026} note, \emph{On the Hodge, Tate, and
standard conjectures for abelian varieties}, is linked from his own
mathematics homepage and available there as
\texttt{math/xnotes/HTS.pdf}. Its Hexagon deposit is
\textbf{2611.00244v1}, submitted 9 November and displayed as published
10 November. Milne states the CM result as Theorem 2.1, p.~4, and
explicitly treats the formerly conditional abelian consequences as
unconditional in §3, p.~6. We read the ten-page note. It is useful
corroboration and context; the numbered primary implications below are
also checked in the original papers.

\subsection{CC-TATE02. Exact abelian consequences of the cited CM
theorem}\label{cc-tate02.-exact-abelian-consequences-of-the-cited-cm-theorem}

Let \(A/\mathbb F_q\) be an abelian variety of arbitrary dimension, let
\(\ell\ne\operatorname{char}\mathbb F_q\), and let
\(N^r(A)_{\mathbb Q}\) denote codimension-\(r\) cycles with rational
coefficients modulo numerical equivalence. The resulting finite-field
statement is

\[
N^r(A)_{\mathbb Q}\otimes_{\mathbb Q}\mathbb Q_\ell
\xrightarrow{\ \sim\ }
H^{2r}_{\mathrm{et}}(A_{\overline{\mathbb F}_q},\mathbb Q_\ell(r))^{G_{\mathbb F_q}}
\qquad\text{for every }r.
\]

In particular, the order of the pole of \(Z(A,t)\) at \(t=q^{-r}\) is
\(\dim_{\mathbb Q}N^r(A)_{\mathbb Q}\). These assertions concern
rational coefficients.

\textbf{Exact original implication.} Milne (1999), \emph{Lefschetz
motives and the Tate conjecture}, \textbf{Theorem 7.1, journal p.~72},
assumes the Hodge conjecture for all complex CM abelian varieties and
proves his Tate statement (0.1) for all abelian varieties over
\(\overline{\mathbb F}_p\). Statement (0.1), \textbf{journal p.~46},
includes both numerical versus homological equivalence and spanning of
the Tate subspace. The following paragraph explicitly gives the
corresponding assertion for each finite-field model and the zeta pole
order. The theorem's proof is organized through Propositions \textbf{7.3
and 7.4, pp.~73--75}. These statements and the complete proofs of these
two propositions were read. The auxiliary theory in §§1--6 is used as
cited literature, not independently reproved. The journal PDF, whose
pagination is 45--76, was inspected; the author's separate erratum
contains an inaccurate bibliographic range 47--81, which is not
propagated here.

\textbf{Finite-field descent check.} The passage from geometric Tate
classes to classes over a given finite field can also be made explicit.
Write an invariant class as a finite \(\mathbb Q_\ell\)-linear
combination of classes of cycles over one common finite extension
\(\mathbb F_{q^m}\). Push each cycle down to \(A\) and divide by \(m\).
On geometric cohomology this is the averaging operator

\[
\frac1m\sum_{j=0}^{m-1}F^j.
\]

It fixes the given \(F\)-invariant class. Thus no enlargement of the
stated ground field is required in the displayed surjectivity assertion.
The numerical-kernel assertion comes from Milne's full Tate statement.
This argument is a formal verification of the field descent, not a new
proof of the Tate input.

\textbf{Standard conjectures.} For every algebraically closed field
\(k\) and every prime \(\ell\ne\operatorname{char}k\), numerical and
\(\ell\)-adic homological equivalence coincide on every abelian variety
over \(k\). For a polarized abelian variety of dimension \(g\), the
primitive rational algebraic classes in codimension \(r\le g/2\) carry
the positive-definite pairing

\[
(a,b)\longmapsto(-1)^r\deg\bigl(a\cdot b\cdot L^{g-2r}\bigr).
\]

The exact inspected locator is Milne (2002), \textbf{Theorem 3.3(a),(b),
pp.~8--9 of the final author manuscript dated 19 September 2001},
\texttt{2002aP.pdf}. Its hypothesis is precisely the CM Hodge theorem in
CC-TATE01. We read the theorem and its full proof, including the
specialization argument to arbitrary algebraically closed \(k\).
\textbf{Corollary 3.4, p.~9}, gives the corresponding polarization of
the category of abelian numerical motives.

For cycles defined over an arbitrary field \(k\), these comparisons
follow by base change to \(\bar k\). Indeed, a complementary cycle over
\(\bar k\) is defined over a finite extension \(k'/k\); the degree
pairing with a cycle defined over \(k\) can be computed using its
pushforward from \(k'\), up to the nonzero extension degree. Thus
numerical triviality over \(k\) is equivalent to numerical triviality
after algebraic closure, with rational coefficients. The \(\ell\)-adic
homological condition is geometric by definition. Positivity on the
rational classes over \(k\) is the restriction of the geometric
positive-definite form. This gives the stated abelian comparison and
positivity over arbitrary fields, for \(\ell\) invertible in the field.

Two source precisions matter. First, van Hoften--Sempliner's §1.1.2
cites \textbf{Theorem 4.6}, the earlier long-version numbering; it
should not be silently represented as the number in the final
manuscript. Second, Milne's abbreviated 2026 display in §3.6 omits the
factor \((-1)^r\); the correct signed form is explicit in the 2002
paper's introduction, p.~2, and §3, p.~6. We use that signed form. The
final manuscript's \textbf{Remark 3.5, p.~9}, concerns an arbitrary
\emph{good} Weil cohomology, where ``good'' means that homological and
numerical equivalence agree on all varieties in the chosen
product-closed class. It does not remove that condition for an
unspecified Weil cohomology theory.

Milne's 2026 note states the finite-field conclusion as \textbf{Theorem
3.1, p.~6}, and the Hodge-standard conclusion as \textbf{Theorem 3.5,
p.~8}. The chain is an application of the cited 2026 theorem to old
conditional theorems. It does not use RH, Dirichlet GRH, a number-field
generalized RH, or an unproved automorphy assertion.

\subsection{CC-TATE03. Rational specialization: a separate 2026
input}\label{cc-tate03.-rational-specialization-a-separate-2026-input}

The Hexagon record \textbf{2610.00226v1}, contributed by \textbf{Keerthi
Madapusi}, is dated \textbf{13 October 2026}. Its actual PDF title is
\emph{Milne's rationality conjecture}, with the byline \textbf{``ChatGPT
and Claude''}; §1.1, p.~3, identifies Madapusi's questions and feedback.
Milne's November note cites this source under Madapusi's name and a
longer descriptive title. The bibliography below records the PDF
attribution and the contributor explicitly.

\textbf{Theorem 1.1, p.~1}, concerns an abelian variety
\(A/\overline{\mathbb Q}\) with good reduction
\(A_0/\overline{\mathbb F}_p\) at a fixed place, a rational Hodge class
\(\gamma\) of codimension \(r\), and a rational Lefschetz class
\(\delta\) of complementary codimension on \(A_0\). The traces

\[
\operatorname{Tr}_v(\gamma_{0,v}\delta_v)
\]

are the images of one rational number for all \(v=\ell\ne p\) and for
the crystalline realization \(v=p\), with the specified coefficient
extension in the latter. Theorem 5.1, pp.~11--12, supplies the
specialized Weil-class dichotomy; §5.1 deduces the main statement. The
lifting input is recorded precisely in \textbf{Proposition A.1,
pp.~13--14}, which invokes Kisin--Zhou, including a quasi-split
replacement and marked tensor transport. These are inspected source
hypotheses and dependency locators; this run does not claim an
independent proof of their lifting arguments.

Combining CC-TATE01 with this manuscript's \textbf{Corollary 6.2,
p.~13}, gives a single cycle \(z_\gamma\in CH^r(A_0)_{\mathbb Q}\)
representing every specialized realization of \(\gamma\). The
corollary's proof was read: CM transport gives
\(\gamma_0=f_0^*\kappa_0\); algebraicity of the CM class \(\kappa\)
permits specialization of a representing cycle; for an actual
homomorphism \(h=Nf_0\), the cycle is \(N^{-2r}h^*Z_0\). This is
algebraicity on the \textbf{special fibre}. It does not assert that
\(\gamma\) is algebraic on a general non-CM characteristic-zero \(A\).

\subsection{CC-TATE04. What this resolves in cote 15, and what it leaves
to
prove}\label{cc-tate04.-what-this-resolves-in-cote-15-and-what-it-leaves-to-prove}

The supplied transcription \texttt{ARCHIVE\_15\_scan100.tex} asks for
sufficient arithmetic conditions for a Frobenius eigenvalue. Scans 102,
104, and 106 then seek a motivic Galois description and its geometric
version. These are distinct mathematical assertions.

\begin{longtable}[]{@{}
  >{\raggedright\arraybackslash}p{(\columnwidth - 4\tabcolsep) * \real{0.3333}}
  >{\raggedright\arraybackslash}p{(\columnwidth - 4\tabcolsep) * \real{0.3333}}
  >{\raggedright\arraybackslash}p{(\columnwidth - 4\tabcolsep) * \real{0.3333}}@{}}
\toprule\noalign{}
\begin{minipage}[b]{\linewidth}\raggedright
Archive issue
\end{minipage} & \begin{minipage}[b]{\linewidth}\raggedright
Consequence supported by the inspected chain
\end{minipage} & \begin{minipage}[b]{\linewidth}\raggedright
Remaining condition
\end{minipage} \\
\midrule\noalign{}
\endhead
\bottomrule\noalign{}
\endlastfoot
Existence of an individual Weil-integer eigenvalue & The Honda--Tate
construction already supplies the occurrence result analyzed separately
in the AQ module. & An occurrence statement does not identify all
correspondences on another variety with that eigenvalue. \\
Tate classes on abelian varieties and their products over finite fields
& The finite-field rational cycle isomorphism in CC-TATE02 holds in
every dimension and codimension. & The 2026 CM theorem is a cited
primary input, not an independently reproved theorem of this report. \\
A smooth projective variety with a motive of abelian type & Milne
(1999), \textbf{Remark 7.6(b), journal p.~76}, explicitly extends Tate
to a variety whose motive \textbf{modulo homological equivalence} lies
in the Tannakian subcategory generated by abelian varieties; products of
curves are an example. & One must establish this actual motivic
membership for the particular variety. Matching Frobenius characteristic
polynomials alone does not establish it. \\
The geometric motivic group in the abelian category & Milne's
\textbf{Proposition 7.3, pp.~73--74}, gives \(M=P\) for the category of
abelian numerical motives over \(\overline{\mathbb F}_p\), with \(P\)
his Weil-number protorus. & This is a statement about the abelian
category and its Tannakian group, not a neutralization of every motivic
gerbe over \(\mathbb Q\). \\
Primitive positivity and the numerical/homological comparison &
CC-TATE02 discharges these conditions for abelian varieties and the
indicated realizations. & No inspected theorem here proves these
statements for an arbitrary smooth projective variety in positive
characteristic. \\
Integral, mixed, or global analytic strengthening & No such
strengthening follows from the rational pure statements just listed. &
Integral comparison, extension data, and number-field analytic
continuation/zeros require their own hypotheses and constructions. \\
\end{longtable}

\textbf{The group-theoretic application in full.} Let \(M\) be the
fundamental group used in §7 of Milne (1999) for abelian numerical
motives over \(\overline{\mathbb F}_p\). The pre-existing construction
gives \(P\subseteq M\subseteq L\). The CM Hodge theorem identifies the
CM Hodge correspondences with algebraic correspondences, so
specialization gives the additional inclusion \(M\subseteq S\) inside
the common Lefschetz group \(T\). Milne's intersection theorem gives
\(S\cap L=P\). Therefore

\[
P\subseteq M\subseteq S\cap L=P,
\]

and hence \(M=P\). This is exactly the known-result application in
Proposition 7.3; it neither assumes nor proves that all smooth
projective motives belong to this category. Milne defines \(P\) by the
Galois module of germs of Weil numbers in §4, \textbf{author-manuscript
p.~14}. These germs use the power transition maps as the finite field
grows. A convention for homological versus cohomological weights must be
retained when comparing his notation with the archive.

\textbf{Gerbe precision.} Milne's separate \textbf{1999b
addendum/erratum, p.~1}, retracts an optimistic introductory claim that
an intersection of Tannakian groups alone automatically constructs the
desired quotient category. Its stated issue is the nonneutral setting.
It does not retract Theorem 7.1; its §7, pp.~3--4, even supplies
Deligne's alternative proof route. Thus the present all-CM theorem
discharges the hypotheses of the inspected abelian Tate theorem, while
an asserted canonical rational fibre functor or the archive's literal
displayed product decomposition must still be justified on its own
terms.

The AQ module makes both sides of this distinction concrete:
\textbf{AQ-ABELIAN-RECOGNITION} proves that full Frobenius determinants
recognize actual rational numerical objects in its explicitly defined
abelian--Artin category, whereas \textbf{AQ-RATIONAL-FIBRE} gives an
elliptic quaternionic obstruction to a rational fibre functor and to
discarding eigenvalue multiplicities. Both complete proofs were
independently read in this run and passed on their stated inputs.

The safe global status statement is therefore precise: \textbf{the
inspected 2026 literature resolves the abelian Tate/standard inputs and
the stated abelian-type Tate input; the unrestricted archive
identification is not established by those inputs.} No assertion that
every broader version is currently open is needed for this conclusion.

\subsection{CC-TATE05. Fourfold confirmation and revision
control}\label{cc-tate05.-fourfold-confirmation-and-revision-control}

Broe's \emph{The Tate conjecture for abelian fourfolds over finite
fields}, \textbf{arXiv:2608.28651v1, 19 August 2026}, gives
\textbf{Theorem 1.1, p.~1}, for every abelian fourfold over a finite
field and \textbf{Corollary 1.2, p.~2}, for standard D on abelian
fourfolds over arbitrary fields, with \(\ell\) invertible. Corollary
8.2, p.~20, repeats the latter with its proof. The complete argument in
§§7--8 was read, with earlier reduction and lifting results retained as
cited inputs. Its \textbf{Remark 1.3, p.~2}, distinguishes positivity on
the numerical quotient from Grothendieck's original homological
formulation. This verifies and strengthens the previously uninspected
Broe entry in SA22; the later all-CM chain makes the abelian dimension
restriction unnecessary for the Tate/standard conclusions used here.

Ningyi Li's separate \textbf{arXiv:2609.27916v2} is explicitly
\textbf{withdrawn on 8 October 2026}. Its abstract must not be used as a
current theorem citation. The withdrawal does not concern the distinct
fivefold manuscript arXiv:2609.06265v2. That manuscript was examined
only selectively and is not required as an input anywhere in this
module.

\subsection{Bibliography and stable primary
URLs}\label{bibliography-and-stable-primary-urls}

Broe, M. (2026, August 19). \emph{The Tate conjecture for abelian
fourfolds over finite fields} (arXiv:2608.28651v1) {[}Preprint{]}.
https://arxiv.org/pdf/2608.28651v1

ChatGPT, \& Claude. (2026, October 13). \emph{Milne's rationality
conjecture} (Hexagon:2610.00226v1) {[}Preprint contributed by Keerthi
Madapusi; authorship and supervision explained in §1.1{]}.
https://hexagonmath.org/pdf/2610.00226v1

Li, N. (2026, October 8). \emph{The Tate conjecture for powers of
abelian fourfolds and the Hodge conjecture for powers of CM fourfolds}
(arXiv:2609.27916v2) {[}Withdrawn preprint record{]}.
https://arxiv.org/abs/2609.27916

Milne, J. S. (1999). Lefschetz motives and the Tate conjecture.
\emph{Compositio Mathematica, 117}(1), 45--76.
https://doi.org/10.1023/A:1000776613765

\begin{itemize}
\tightlist
\item
  Inspected journal PDF:
  https://www.cambridge.org/core/services/aop-cambridge-core/content/view/962D833FAAC3219F55A58B193323B52A/S0010437X99000317a.pdf/lefschetz-motives-and-the-tate-conjecture.pdf
\item
  Author manuscript: https://jmilne.org/math/articles/1999bP.pdf
\item
  Author addendum/erratum:
  https://jmilne.org/math/articles/add/1999b.pdf
\end{itemize}

Milne, J. S. (2002). Polarizations and Grothendieck's standard
conjectures. \emph{Annals of Mathematics, 155}(2), 599--610.
https://doi.org/10.2307/3062126

\begin{itemize}
\tightlist
\item
  Inspected final author manuscript, dated September 19, 2001:
  https://www.jmilne.org/math/articles/2002aP.pdf
\item
  Publisher metadata: https://annals.math.princeton.edu/2002/155-2/p07
\item
  Author's separate explanatory addendum:
  https://www.jmilne.org/math/articles/2002a.html
\end{itemize}

Milne, J. S. (2026, November 9). \emph{On the Hodge, Tate, and standard
conjectures for abelian varieties} (Hexagon:2611.00244v1) {[}Expository
preprint{]}. https://www.jmilne.org/math/xnotes/HTS.pdf

\begin{itemize}
\tightlist
\item
  Versioned deposit: https://hexagonmath.org/2611.00244v1
\end{itemize}

van Hoften, P., \& Sempliner, J. (2026, October 18). \emph{The Hodge
conjecture for CM abelian varieties} (Hexagon:2610.00259v2)
{[}Preprint{]}. https://hexagonmath.org/pdf/2610.00259v2

The source snapshots, extracted text, and selected rendered theorem
pages are preserved in \texttt{research/cc\_sources/}. Exact
web-reference identifiers and inspection methods are in the accompanying
citation and check-event ledgers.
